\chapter{Lagrangian Fibrations and Integral Affine Manifolds}
\section{Symplectic Geometry}

In this section we will recall some of the basics of symplectic geometry and
Lagrangian submanifolds.

\begin{definition}[Symplectic Manifold]
    A symplectic manifold $(M, \omega)$ consists of a smooth manifold $M$ and a closed 2-form $\omega$ such that $\omega$ is non-degenerate.
\end{definition}

Symplectic geometry has its roots in classical mechanics as the phase space of
a dynamical system naturally has the structure of a symplectic manifold. TODO:
elaborate more on symplectic geometry - Hamiltons equations, applications to
topology, Poisson geometry.

Non-degeneracy of the symplectic form gives a canonical way to identify the
tangent and cotangent bundle of a symplectic manifold. This is done via the map
$$ \tilde{\omega}|_p : T_pM \to T_p^*M $$ $$ v \mapsto \omega_p(v, \cdot)$$ for
$v \in T_pM$. This is an isomorphism by non-degeneracy and so, namely, it is
invertible. This identification leads to the notion of a Hamiltonian vector
field.

\begin{definition}[Hamiltonian Vector Field]
    For any $H\in C^{\infty}(M)$ we can associate a unique Hamiltonian vector field by
    $$\iota_{X_H}\omega := \omega(X_H, \cdot) = dH$$
\end{definition}

For mathematicians there is generally no distinguished choice of Hamiltonian
function $H$. However, physicists often pick a Hamiltonian function based on
physical criteria and are generally interested in studying the flow of this
Hamiltonian vector field which, for instance, recovers Hamilton's equations.

\subsection{The Cotangent Bundle}
One of the most important examples of symplectic manifolds is the cotangent
bundle, which naturally inherits the structure of a symplectic manifold. In
fact, Darbouxs theorem states that every symplectic manifold looks locally like
a cotangent bundle. This class of symplectic manifolds will playu a big role
throughout this thesis and we will see now exactly how the cotangent bundle
obtains a symplectic structure.

\begin{eg}[Cotangent Bundle]
    Let $Q$ be any smooth manifold and $M = T^*Q$ its cotangent bundle.
    We will now illustrate that there is a canonical one form $\theta \in \Omega^1(T^*Q)$
    and show that the exterior derivative $-d\theta$ defines a symplectic structure on $T^*Q$.

    Let $\pi : T^*Q \to Q$ be the canonical projection map. Then for $X_p \in
        T_p(T^*Q)$ we define $\langle \theta_p , X_p \rangle := \langle p , d_p\pi(X_p)
        \rangle$ where $\langle \cdot , \cdot \rangle$ denotes contraction of a vector
    and covector. This canonical one-form satisfies the following important
    property.
    \begin{lemma}
        $\theta$ is the unique one-form such that for any other one-form $\alpha \in \Omega^1(Q)$
        $$\alpha = \alpha^* \theta$$
        where we view $\alpha$ as a section $\alpha: Q \to T^*Q$.
    \end{lemma}
    \begin{proof}
        Let $w \in T_qQ$. Then we have that
        \begin{align*}
            \langle (\alpha^*\theta)_q , w \rangle & = \langle \theta_{\alpha_q} , d_q\alpha (w) \rangle \\
                                                   & = \langle \alpha_q , w \rangle
        \end{align*}
        Uniqueness is obtained by noting that every tangent vector $v \in T_p(T^*Q)$ can be written in the form $d_q\alpha(w)$
        except for those in the kernel of $d_p\pi$. But the vectors not in this kernel form an open, dense set and so by continuity we get uniqueness.
    \end{proof}

    We can also find an expression for this canonical one-form in coordinates.

    \begin{lemma}
        In local cotangent coordinates $q_1, \dots , q_n, p_1, \dots , p_n$ on $T^*Q$, the canonical one-form is given by
        $$\theta|_{T^*U} = \sum_j p_j dq_j$$
    \end{lemma}
    \begin{proof}
        We let $\alpha = \sum_j \alpha_j dq_j$ be a $1$-form on $U$. We have that $\alpha^*p_j = \alpha_j$ and $\alpha^*q_j = q_j$ which gives that
        \begin{align*}
            \alpha^*\theta|_{T^*U} & = \alpha^* \sum_j p_j dq_j \\
                                   & =\sum_j \alpha_j dq_j      \\
                                   & = \alpha
        \end{align*}
    \end{proof}

    Putting this together we obtain a symplectic form on $T^*Q$.
    \begin{theorem}
        The $2$-form $-d\theta$ defines a symplectic form on $T^*Q$.
    \end{theorem}
    \begin{proof}
        In local coordinates we have that $\omega = \sum_j dq_j \wedge dp_j$ which is clearly closed and non-degenerate.
    \end{proof}
\end{eg}

We will also wish to speak about maps between the symplectic manifolds. The
correct notion of "isomorphism" in this category is the following.

\begin{definition}
    Let $(M, \omega)$ and $(N, \alpha)$ be two symplectic manifolds. Then a symplecto-morphism between $M$ and $N$ is a diffeomorphism
    $$\varphi: M \to N$$
    such that $\omega = \varphi^* \alpha$.
\end{definition}

\subsection{Lagrangian Submanifolds}
The goal of this thesis is to classify Lagrangian fibrations over compact
surfaces. In this section we will introduce the concept of Lagrangian
submanifolds and prove some basic results regarding these.
\begin{definition}[Symplectic Orthogonal]
    Let $(M, \omega)$ be a symplectic manifold and let $N \subset M$ be a submanifold. The symplectic orthogonal at a point $p \in M$ is
    $$(T_pN)^{\omega} := \{v \in T_pM \mid \omega(v, w) = 0, \forall w \in N\}$$
\end{definition}

\begin{itemize}
    \item If $(T_pN)^{\omega} \subset T_pN$ then $T_pN$ is said to be
          \emph{co-isotropic}.
    \item If $T_pN\subset (T_pN)^{\omega} $ then $T_pN$ is said to be \emph{isotropic}.
    \item If $(T_pN)^{\omega} = T_pN$ then $T_pN$ is said to be \emph{Lagrangian}.
    \item If $(T_pN)^{\omega} \cap T_pN = \empty$ then $T_pN$ is said to be
          \emph{symplectic}.
\end{itemize}

A submanifold $N$ of a symplectic manifold $M$ is said to be co-isotropic
(isotropic, Lagrangian, symplectic) if its tangent space is a co-isotropic
(isotropic, Lagrangian, symplectic) subspace of $T_pM$ at every point $p \in
    M$.

\begin{remark}
    We note that a submanifold $\iota : N \to M$ is isotropic if $\iota^*\omega =0$ and furthermore, it is Lagrangian if we also have
    $\text{dim}N = \frac{1}{2}\text{dim}N$.
\end{remark}

We will make use of the following lemma.
\begin{lemma}
    \label{lemma:cotangent-lagrangian-sub}
    The graph $\Gamma_{\alpha} \subset T^*Q$ of a 1-form $\alpha$ is Lagrangian if and only if $\alpha$ is closed.
\end{lemma}
\begin{proof}
    It is clear that $\text{dim}\Gamma_{\alpha} =\frac{1}{2}\text{dim}T^*Q$.
    \begin{align*}
        \alpha^*\omega & = -\alpha^*d\theta   \\
                       & = - d\alpha^* \theta \\
                       & = - d\alpha
    \end{align*}
    and so $\alpha^* \omega = 0 $ if and only if $\alpha$ is closed completing the proof by the previous remark.
\end{proof}

The cotangent bundle $T^*Q$ also gives rise to another Lagrangian submanifold.

\begin{eg}
    Let $\pi : T^*Q \to Q$ and fix $x \in Q$. Then $\pi^{-1}(x) = T_x^*Q$ is a Lagrangian submanifold inside $(T^*Q, \omega)$. It is easy to see that $\text{dim }T_x^*Q = \frac{1}{2}\text{dim }T^*Q $.
    Letting $i: T_x^*Q \hookrightarrow T^*Q$ be the inclusion and considering local coordinates $(q_1, \dots , q_n, p_1 \dots p_n)$ around $x$ and writing $\omega = \sum_j dq_j \wedge dp_j$ it follows that $\omega|_{T_x^*Q} = 0$ as $i^*(\sum_j dq_j \wedge dp_j) = \sum_j d(i^*q_j) \wedge d(i^*p_j)$ and
    $d(i^*q_i) = 0$ as $i^*q_i$ is constant.
\end{eg}

\section{Lagrangian Fibrations}

One of the goals of this thesis is to classify the Lagrangian fibrations over
closed surfaces. In this section we will give an overview of these objects and
give some examples and basic results surrounding them. We will then describe
how the base space of a Lagrangian fibration inherits an integral affine
structure.

\subsection{Definition of Lagrangian Fibration}
\begin{definition}
    A Lagrangian fibration $\pi: (M, \omega) \to B $ is a surjective map whose regular fibres are Lagrangian
    submanifolds of $(M, \omega)$.
\end{definition}

\begin{eg}
    We seen in Lemma \ref{lemma:cotangent-lagrangian-sub} that the fibres of $\pi: (T^*Q, \omega_{can}) \to Q$ are Lagrangian submanifolds
    and hence this is a Lagrangian fibration.
\end{eg}

\begin{eg}
    Let $Q= (\R / \Z)^n = \mathbb{T}^n$ be the n-torus. Then we have that $T^*(\mathbb{T}^n) = \mathbb{T}^n \times \R^n$ and the map
    $$
        \pi : T^*(\mathbb{T}^n) \to \R^n
    $$
    is a Lagrangian fibration with fibre $\mathbb{T}^n$.
\end{eg}
\begin{definition}[Shift by a 1-form.]
    Let $\pi : (T^*Q, \omega_{can}) \to Q$ and fix $\alpha \in \Omega^1(Q)$. We define the \emph{shift by $\alpha$} to be the map
    $$ s_{\alpha}: T^*Q \to T^*Q$$
    $$ \beta_q \mapsto \beta_q + \alpha_q$$
\end{definition}
\begin{prop}
    This defines a symplecto-morphism of $T^*Q$ if and only if $d\alpha = 0$.
\end{prop}
\begin{proof}
    This map is a fibrewise translation and hence is obviously a diffeomorphism. We now check when it preserves
    the symplectic form. We calculate
    \begin{align*}
        s_{\alpha}^*\omega_{can} & = -s_{\alpha}^*d\theta  \\
                                 & = -ds_{\alpha}^* \theta \\
    \end{align*}
    We calculate
    \begin{align*}
        s_{\alpha}^* \theta (v)|_{\beta_q} & = \theta_{\beta_q + \alpha_q}(ds_{\alpha}(v))     \\
                                           & = (\beta_q +\alpha_q)(d(\pi \circ s_{\alpha})(v)) \\
                                           & = (\beta_q +\alpha_q)(v)                          \\
                                           & = \theta_{\beta_q}(v) + \alpha_q(v)
    \end{align*}
    and so $s_{\alpha}^* \theta = \theta + \alpha$. Now
    \begin{align*}
        -ds_{\alpha}^* \theta & = -d\theta - d\alpha     \\
                              & = \omega_{can} - d\alpha
    \end{align*}
    and hence this defines a symplecto-morphism if and only if $d\alpha = 0$.
\end{proof}
\subsection{Integral Affine Structure on the Base}

\subsection{Action Angle Coordinates}